% !TeX program = platex
% Overleaf: Compiler = LaTeX / Main document = main.tex
\documentclass[dvipdfmx,a4paper,11pt]{jsarticle}
\usepackage[margin=23mm]{geometry}
\usepackage{booktabs}
\usepackage{fancyvrb}
\usepackage[dvipdfmx,unicode,colorlinks=true,urlcolor=blue,linkcolor=black,citecolor=black]{hyperref}
\usepackage{pxjahyper}
\fvset{fontsize=\scriptsize,tabsize=4}
\newcommand{\code}[1]{\texttt{#1}}
\title{デザインパターンを使ったディズニー作品の鑑賞リスト管理}
\author{学籍番号：5424029{\hspace{35mm}}\quad 氏名：斎藤柾宏{\hspace{35mm}}}
\date{\today}
\begin{document}
\maketitle
\section{はじめに}
映画の鑑賞リストを作るとき、作品名をただ並べるだけでなく、シリーズごとにまとめたり、見終わった作品の数を確認したりできると便利だと考えた。
例えば「アナと雪の女王」と「アナと雪の女王2」を同じシリーズとしてまとめ、何作品見終わったかを表示できるようにする。
そこで、ディズニーのアニメーション作品を題材として、作品の分類と鑑賞状況を管理するプログラムを作成した。

作品とグループを同じように扱う部分にはCompositeパターンを使い、未鑑賞・鑑賞中・鑑賞済みという状態ごとの処理にはStateパターンを使った。
さらに、一度見終わった作品をもう一度見る場合も考え、再鑑賞を始めたら状態が鑑賞中に戻るようにした。

本レポートでは、条件分岐を使ったコードと、2つのパターンを使ったコードを比較し、
同じ作品を登録し、同じ順番で操作して結果を確かめたうえで、コードの違いや機能追加のしやすさについて考察することを目的としている。

\section{背景}
映画やアニメの作品数が増えると、どの作品を見たのか、まだ見ていない作品は何かを把握しにくくなることがある。特にディズニー作品には、「アナと雪の女王」や「シュガー・ラッシュ」のように続編がある作品もある。そのため、作品を一つずつ管理するだけでなく、シリーズごとにまとめて鑑賞状況を確認できると便利だと考えた。

ただ、作品とシリーズを一つのクラスで扱おうとすると、それが作品なのかグループなのかを毎回判定する必要がある。また、未鑑賞・鑑賞中・鑑賞済みでは実行できる操作が異なるため、状態ごとに条件分岐を書くことになる。作品の分類や鑑賞状態が増えた場合には、こうした処理が分かりにくくなる可能性があると考えた。

そこで今回は、作品とグループの管理にCompositeパターンを、鑑賞状態による処理の切り替えにStateパターンを使うことにした。パターンを使わない場合と使う場合、両方を比較することで、どの部分が整理されるのか、反対に実装が複雑になる部分はないかを確かめることとする。

\section{作成する鑑賞リスト}
\subsection{今回作成した鑑賞リストの分類}
実行例では、全体リストを「近年のアニメーション作品」と「名作アニメーション作品」に分け、計12作品を登録することにした。
前者には2つのシリーズと4つの単独作品、後者には4作品を配置した。
登録した作品の構成を以下に示す。\newpage
\begin{Verbatim}
ディズニー作品の鑑賞リスト
  近年のアニメーション作品
    アナと雪の女王シリーズ
      アナと雪の女王
      アナと雪の女王2
    シュガー・ラッシュシリーズ
      シュガー・ラッシュ
      シュガー・ラッシュ：オンライン
    塔の上のラプンツェル
    ズートピア
    ベイマックス
    モアナと伝説の海
  名作アニメーション作品
    ライオン・キング
    アラジン
    美女と野獣
    リトル・マーメイド
\end{Verbatim}
\subsection{鑑賞状況の管理}
初期状態は未鑑賞とし、表\ref{tab:state}の操作を行う。
それ以外の操作では拒否メッセージを表示し、状態を変えないことにした。
\begin{table}[htbp]\centering
\caption{使用する状態遷移}\label{tab:state}
\begin{tabular}{lll}\toprule
現在の状態 & 操作 & 遷移先\\\midrule
未鑑賞 & \code{start()}：開始 & 鑑賞中\\
鑑賞中 & \code{finish()}：完了 & 鑑賞済み\\
鑑賞済み & \code{rewatch()}：再鑑賞 & 鑑賞中\\\bottomrule
\end{tabular}
\end{table}\\
集計には現在の状態を使う。つまり再鑑賞を始めた作品は鑑賞中に戻り、鑑賞済み件数から外れる。

\section{プログラム仕様}
ディズニーのアニメーション作品12本を登録し、分類やシリーズごとにまとめて管理する。
各作品は「未鑑賞」「鑑賞中」「鑑賞済み」のいずれかの状態を持ち、鑑賞の開始・完了・再鑑賞を行える。
再鑑賞を始めた作品は鑑賞中に戻り、その状態では行えない操作を指定した場合はメッセージを表示して状態を変えない。

一覧には各作品の状態と、グループごとの鑑賞済み作品数を表示する。
グループの中に別のグループを登録することも可能とし、再鑑賞中の作品は鑑賞済み件数に含めない。
操作はプログラム内で決めた順序で実行し、パターン適用前後で同じ結果になるかを比較する。
キーボードからの入力は行わず、登録したリストや操作結果をコンソールに表示する。
鑑賞操作は個々の作品に対して行い、シリーズ全体を一度に鑑賞済みにする機能は設けない。

\section{実行結果}
JDK 21.0.11で両版をコンパイルし、同じリスト構成と操作順で実行した。
実行したところ、同じ結果となった。以下に実行結果を示す。
\begin{Verbatim}
[ディズニー作品の鑑賞リスト] 0/12作品鑑賞済み           
  [近年のアニメーション作品] 0/8作品鑑賞済み
    [アナと雪の女王シリーズ] 0/2作品鑑賞済み
      アナと雪の女王：未鑑賞
      アナと雪の女王2：未鑑賞
    [シュガー・ラッシュシリーズ] 0/2作品鑑賞済み
      シュガー・ラッシュ：未鑑賞
      シュガー・ラッシュ：オンライン：未鑑賞
    塔の上のラプンツェル：未鑑賞
    ズートピア：未鑑賞
    ベイマックス：未鑑賞
    モアナと伝説の海：未鑑賞
  [名作アニメーション作品] 0/4作品鑑賞済み
    ライオン・キング：未鑑賞
    アラジン：未鑑賞
    美女と野獣：未鑑賞
    リトル・マーメイド：未鑑賞
--- 鑑賞操作 ---
アナと雪の女王：完了不可（未鑑賞）
アナと雪の女王：再鑑賞不可（未鑑賞）
アナと雪の女王：鑑賞を開始しました
アナと雪の女王：開始不可（鑑賞中）
アナと雪の女王：鑑賞を完了しました
アナと雪の女王：完了不可（鑑賞済み）
アナと雪の女王：開始不可（鑑賞済み）
アナと雪の女王：再鑑賞を開始しました
[ディズニー作品の鑑賞リスト] 0/12作品鑑賞済み
  [近年のアニメーション作品] 0/8作品鑑賞済み
    [アナと雪の女王シリーズ] 0/2作品鑑賞済み
      アナと雪の女王：鑑賞中
      アナと雪の女王2：未鑑賞
    [シュガー・ラッシュシリーズ] 0/2作品鑑賞済み
      シュガー・ラッシュ：未鑑賞
      シュガー・ラッシュ：オンライン：未鑑賞
    塔の上のラプンツェル：未鑑賞
    ズートピア：未鑑賞
    ベイマックス：未鑑賞
    モアナと伝説の海：未鑑賞
  [名作アニメーション作品] 0/4作品鑑賞済み
    ライオン・キング：未鑑賞
    アラジン：未鑑賞
    美女と野獣：未鑑賞
    リトル・マーメイド：未鑑賞
アナと雪の女王：鑑賞を完了しました
アナと雪の女王2：鑑賞を開始しました
塔の上のラプンツェル：鑑賞を開始しました
塔の上のラプンツェル：鑑賞を完了しました
シュガー・ラッシュ：鑑賞を開始しました
シュガー・ラッシュ：鑑賞を完了しました
ライオン・キング：鑑賞を開始しました
ライオン・キング：鑑賞を完了しました
ズートピア：鑑賞を開始しました
--- 最終一覧 ---
[ディズニー作品の鑑賞リスト] 4/12作品鑑賞済み
  [近年のアニメーション作品] 3/8作品鑑賞済み
    [アナと雪の女王シリーズ] 1/2作品鑑賞済み
      アナと雪の女王：鑑賞済み
      アナと雪の女王2：鑑賞中
    [シュガー・ラッシュシリーズ] 1/2作品鑑賞済み
      シュガー・ラッシュ：鑑賞済み
      シュガー・ラッシュ：オンライン：未鑑賞
    塔の上のラプンツェル：鑑賞済み
    ズートピア：鑑賞中
    ベイマックス：未鑑賞
    モアナと伝説の海：未鑑賞
  [名作アニメーション作品] 1/4作品鑑賞済み
    ライオン・キング：鑑賞済み
    アラジン：未鑑賞
    美女と野獣：未鑑賞
    リトル・マーメイド：未鑑賞
\end{Verbatim}

\section{比較}
パターン適用前は、作品とグループの管理や状態の処理が1つのクラスにまとまり、コード全体を追いやすいところがある。
クラス数も少ないため、機能が少なく変更がない場合は簡単に実装できる。
一方、作品とグループの判定や状態の条件分岐が複数箇所にあり、規則を変更するときはそれぞれ確認する必要があるというデメリットが存在した。

適用後は、Compositeによって作品とグループを同じ方法で表示・集計できるようになった。また、Stateによって状態ごとの処理がまとまり、確認や修正がしやすくなった。クラス数は増えるが、今後、作品の種類や状態を増やすときに、変更する場所を分かりやすくする効果がある。

\newpage
\section{考察}
今回は、作品をグループに分ける処理と、鑑賞状態を管理する処理に2つのパターンを使った。
作品を見終えたり再鑑賞を始めたりすると、グループの鑑賞済み件数にも反映される。
そのため、シリーズの中で何作品を見終えたかを確認しやすい。
ただし、状態が3種類だけなら、パターンを使わず条件分岐で書く方法でも対応できる。

パターンを使うとクラスが増えるため、今後増やしたい機能に合わせて使うかどうかを考える必要がある。
また、一度見た作品を記録しておきたい場合は、現在の状態とは別に履歴を残す必要がある。
実際に使うには、同じ作品を重ねて登録したり、グループの中身がループしたりしないようにする工夫も必要である。
今後は、空のグループや別の操作順でも正しく動くかを確認したい。

\section{まとめ}
今回は、CompositeとStateを使ってディズニー作品の鑑賞リストを管理するプログラムを作成した。
適用前後で同じ結果となり、作品の分類と鑑賞状態の処理を分けて整理できた。
クラス数は増えるが、機能を追加するときに変更する場所が分かりやすくなると考える。

\begin{thebibliography}{9}
\bibitem{1} 尾崎知伸『デザインパターンの応用』\hspace{1zw}最終閲覧日 : \date{\today}\\\url{https://nu.instructure.com/courses/103710/pages/%E3%83%87%E3%82%B6%E3%82%A4%E3%83%B3%E3%83%91%E3%82%BF%E3%83%BC%E3%83%B3%E8%B3%87%E6%96%99?module_item_id=520525}
\bibitem{2} 尾崎知伸『デザインパターンの実践』\hspace{1zw}最終閲覧日 : \date{\today}\\\url{https://nu.instructure.com/courses/103710/pages/%E3%83%87%E3%82%B6%E3%82%A4%E3%83%B3%E3%83%91%E3%82%BF%E3%83%BC%E3%83%B3%E8%B3%87%E6%96%99?module_item_id=520525}
\end{thebibliography}

\appendix
\clearpage
\section{パターン適用前のコード全文}
\begin{Verbatim}
package before;

import java.util.ArrayList;
import java.util.List;

/** 適用前：種類フラグと状態の条件分岐によって鑑賞リストを管理する。 */
public class BeforeMain {
    enum Status { UNWATCHED, WATCHING, WATCHED }

    static class Entry {
        private final String title;
        private final boolean group;
        private final List<Entry> children = new ArrayList<>();
        private Status status = Status.UNWATCHED;

        Entry(String title, boolean group) {
            this.title = title;
            this.group = group;
        }

        void add(Entry child) {
            if (!group) throw new IllegalStateException("作品には追加できません");
            children.add(java.util.Objects.requireNonNull(child));
        }

        int totalCount() {
            if (!group) return 1;
            int count = 0;
            for (Entry child : children) count += child.totalCount();
            return count;
        }

        int watchedCount() {
            if (!group) return status == Status.WATCHED ? 1 : 0;
            int count = 0;
            for (Entry child : children) count += child.watchedCount();
            return count;
        }

        String label() {
            switch (status) {
                case UNWATCHED: return "未鑑賞";
                case WATCHING: return "鑑賞中";
                case WATCHED: return "鑑賞済み";
                default: throw new IllegalStateException();
            }
        }

        void print(String indent) {
            if (group) {
                System.out.printf("%s[%s] %d/%d作品鑑賞済み%n",
                        indent, title, watchedCount(), totalCount());
                for (Entry child : children) child.print(indent + "  ");
            } else {
                System.out.println(indent + title + "：" + label());
            }
        }

        void start() {
            requireMovie();
            if (status == Status.UNWATCHED) {
                status = Status.WATCHING;
                message("鑑賞を開始しました");
            } else reject("開始");
        }

        void finish() {
            requireMovie();
            if (status == Status.WATCHING) {
                status = Status.WATCHED;
                message("鑑賞を完了しました");
            } else reject("完了");
        }

        void rewatch() {
            requireMovie();
            if (status == Status.WATCHED) {
                status = Status.WATCHING;
                message("再鑑賞を開始しました");
            } else reject("再鑑賞");
        }

        private void requireMovie() {
            if (group) throw new IllegalStateException("作品を選択してください");
        }
        private void message(String text) { System.out.println(title + "：" + text); }
        private void reject(String action) { message(action + "不可（" + label() + "）"); }
    }

    public static void main(String[] args) {
        Entry root = new Entry("ディズニー作品の鑑賞リスト", true);
        Entry modern = new Entry("近年のアニメーション作品", true);
        Entry series = new Entry("アナと雪の女王シリーズ", true);
        Entry ralphSeries = new Entry("シュガー・ラッシュシリーズ", true);
        Entry classics = new Entry("名作アニメーション作品", true);
        Entry first = new Entry("アナと雪の女王", false);
        Entry second = new Entry("アナと雪の女王2", false);
        Entry ralph = new Entry("シュガー・ラッシュ", false);
        Entry ralph2 = new Entry("シュガー・ラッシュ：オンライン", false);
        Entry rapunzel = new Entry("塔の上のラプンツェル", false);
        Entry zootopia = new Entry("ズートピア", false);
        Entry baymax = new Entry("ベイマックス", false);
        Entry moana = new Entry("モアナと伝説の海", false);
        Entry lion = new Entry("ライオン・キング", false);
        Entry aladdin = new Entry("アラジン", false);
        Entry beauty = new Entry("美女と野獣", false);
        Entry mermaid = new Entry("リトル・マーメイド", false);
        root.add(modern);
        root.add(classics);
        modern.add(series);
        modern.add(ralphSeries);
        modern.add(rapunzel);
        modern.add(zootopia);
        modern.add(baymax);
        modern.add(moana);
        series.add(first);
        series.add(second);
        ralphSeries.add(ralph);
        ralphSeries.add(ralph2);
        classics.add(lion);
        classics.add(aladdin);
        classics.add(beauty);
        classics.add(mermaid);
        root.print("");
        System.out.println("--- 鑑賞操作 ---");
        first.finish();
        first.rewatch();
        first.start();
        first.start();
        first.finish();
        first.finish();
        first.start();
        first.rewatch();
        root.print("");
        first.finish();
        second.start();
        rapunzel.start();
        rapunzel.finish();
        ralph.start();
        ralph.finish();
        lion.start();
        lion.finish();
        zootopia.start();
        System.out.println("--- 最終一覧 ---");
        root.print("");
    }
}
\end{Verbatim}
\clearpage
\section{パターン適用後のコード全文}
\begin{Verbatim}
package after;

import java.util.ArrayList;
import java.util.List;

/** 適用後：Compositeで階層を、Stateで作品の鑑賞状態を管理する。 */
public class AfterMain {
    // Component：作品とグループに共通の操作。
    interface WatchEntry {
        int totalCount();
        int watchedCount();
        void print(String indent);
    }

    // Composite：子の種類を判定せず、表示と集計を再帰的に行う。
    static class WatchGroup implements WatchEntry {
        private final String title;
        private final List<WatchEntry> children = new ArrayList<>();

        WatchGroup(String title) { this.title = title; }
        void add(WatchEntry child) {
            children.add(java.util.Objects.requireNonNull(child));
        }
        public int totalCount() {
            int count = 0;
            for (WatchEntry child : children) count += child.totalCount();
            return count;
        }
        public int watchedCount() {
            int count = 0;
            for (WatchEntry child : children) count += child.watchedCount();
            return count;
        }
        public void print(String indent) {
            System.out.printf("%s[%s] %d/%d作品鑑賞済み%n",
                    indent, title, watchedCount(), totalCount());
            for (WatchEntry child : children) child.print(indent + "  ");
        }
    }

    // State：不許可操作の共通処理。許可する操作を各状態で上書きする。
    interface ViewingState {
        String label();
        default int watchedCount() { return 0; }
        default void start(Movie movie) { movie.reject("開始", label()); }
        default void finish(Movie movie) { movie.reject("完了", label()); }
        default void rewatch(Movie movie) { movie.reject("再鑑賞", label()); }
    }

    static class Unwatched implements ViewingState {
        public String label() { return "未鑑賞"; }
        public void start(Movie movie) {
            movie.changeState(new Watching());
            movie.message("鑑賞を開始しました");
        }
    }

    static class Watching implements ViewingState {
        public String label() { return "鑑賞中"; }
        public void finish(Movie movie) {
            movie.changeState(new Watched());
            movie.message("鑑賞を完了しました");
        }
    }

    static class Watched implements ViewingState {
        public String label() { return "鑑賞済み"; }
        public int watchedCount() { return 1; }
        public void rewatch(Movie movie) {
            movie.changeState(new Watching());
            movie.message("再鑑賞を開始しました");
        }
    }

    // CompositeのLeafであり、StateのContextでもある。
    static class Movie implements WatchEntry {
        private final String title;
        private ViewingState state = new Unwatched();

        Movie(String title) { this.title = title; }
        private void changeState(ViewingState next) { state = next; }
        private void message(String text) { System.out.println(title + "：" + text); }
        private void reject(String action, String label) {
            message(action + "不可（" + label + "）");
        }
        void start() { state.start(this); }
        void finish() { state.finish(this); }
        void rewatch() { state.rewatch(this); }
        public int totalCount() { return 1; }
        public int watchedCount() { return state.watchedCount(); }
        public void print(String indent) {
            System.out.println(indent + title + "：" + state.label());
        }
    }

    public static void main(String[] args) {
        WatchGroup root = new WatchGroup("ディズニー作品の鑑賞リスト");
        WatchGroup modern = new WatchGroup("近年のアニメーション作品");
        WatchGroup series = new WatchGroup("アナと雪の女王シリーズ");
        WatchGroup ralphSeries = new WatchGroup("シュガー・ラッシュシリーズ");
        WatchGroup classics = new WatchGroup("名作アニメーション作品");
        Movie first = new Movie("アナと雪の女王");
        Movie second = new Movie("アナと雪の女王2");
        Movie ralph = new Movie("シュガー・ラッシュ");
        Movie ralph2 = new Movie("シュガー・ラッシュ：オンライン");
        Movie rapunzel = new Movie("塔の上のラプンツェル");
        Movie zootopia = new Movie("ズートピア");
        Movie baymax = new Movie("ベイマックス");
        Movie moana = new Movie("モアナと伝説の海");
        Movie lion = new Movie("ライオン・キング");
        Movie aladdin = new Movie("アラジン");
        Movie beauty = new Movie("美女と野獣");
        Movie mermaid = new Movie("リトル・マーメイド");
        root.add(modern);
        root.add(classics);
        modern.add(series);
        modern.add(ralphSeries);
        modern.add(rapunzel);
        modern.add(zootopia);
        modern.add(baymax);
        modern.add(moana);
        series.add(first);
        series.add(second);
        ralphSeries.add(ralph);
        ralphSeries.add(ralph2);
        classics.add(lion);
        classics.add(aladdin);
        classics.add(beauty);
        classics.add(mermaid);
        root.print("");
        System.out.println("--- 鑑賞操作 ---");
        first.finish();
        first.rewatch();
        first.start();
        first.start();
        first.finish();
        first.finish();
        first.start();
        first.rewatch();
        root.print("");
        first.finish();
        second.start();
        rapunzel.start();
        rapunzel.finish();
        ralph.start();
        ralph.finish();
        lion.start();
        lion.finish();
        zootopia.start();
        System.out.println("--- 最終一覧 ---");
        root.print("");
    }
}
\end{Verbatim}
\clearpage
\section{生成AIの利用に関する記録}
\subsection{利用範囲と記録の状態}
本課題では、OpenAIの対話型AI（Codex）を、テーマの検討、授業資料の確認、
Javaコードと説明資料の生成、実行確認、LaTeXレポートの作成に利用した。
本文とコードは生成AIが作成した原稿を含む。
下記は作成時点で確認できる会話の要約であり、完全な会話記録ではない。
生成AIを使わずに本人だけで作成した旨を主張するものではない。

\textbf{提出前の作業：}課題資料は生成AIの利用履歴を末尾に添付することを求めている。
この要約のみでその条件を満たすとは判断せず、今回のレポート作成依頼と応答も含む実際の会話履歴を
原文のまま本付録の末尾へ追記する必要がある。
IDEが自動で付加する情報や長いツール出力を省略する場合も、省略箇所と範囲を明記し、履歴を改変して見せないこと。

\subsection{会話経過の要約（原文ではない）}
\begin{enumerate}
\item 利用者は課題PDFに基づくテーマの提案を依頼した。AIは最初、授業資料の範囲を確認せずStrategyやObserverを候補に含めた。
\item 利用者は、パターンを同じ授業フォルダのPDFにあるものに限定するよう指定した。
AIは授業資料を確認し、先の候補の不適合を訂正し、DecoratorとStateを使う模擬店注文管理を提案した。
\item 利用者は学園祭の模擬店注文管理を選び、Javaで適用前後のコードをコードフォルダへ作るよう依頼した。
AIは2つのJavaファイルとREADMEを作成し、コンパイルと実行を行った。
\item 動作確認の初回では検証用の想定行数を15行としていたが、実際は13行だった。
その後のUTF-8指定ではPowerShellの引数解釈による失敗もあり、引数を引用して再実行した。
最終的に両版が正常終了し、13行の出力が一致することを確認した。
\item 利用者はbuildフォルダの意味と実行方法を尋ね、AIはclassファイルの生成先およびコンパイル・実行コマンドを説明した。
\item 利用者は採点基準への適合を尋ねた。AIは基本的なコード条件を満たす一方、独自性、本人の考察、PDFと履歴の準備が必要であると説明した。
また、同一トッピングを重ねる機能の前後差を明示した。
\item 利用者は「レポートをOverLeaf形式で作成して(Overleafフォルダに)」と依頼し、AIが本LaTeX原稿を作成した。
\item 利用者はコンパイルの問題を報告し、2ファイル構成と軽量化を依頼した。AIはコードを本文へ統合し、upLaTeXと簡素なコード表示へ変更した。PDFの組版はローカルでは未確認である。
\item 利用者はディズニー作品の鑑賞リスト管理へ題材を変更し、AIがCompositeとStateを用いる前後のJavaコードを生成した。実行例では出力27行の一致を確認した。
\item 利用者は変更後のコードに合わせたmain.texの書換えを依頼し、AIが本文・比較・実行結果・コード全文を更新した。
\item 利用者は鑑賞リストの拡充を依頼し、AIは12作品と複数階層へ拡張してJavaコード・本文・実行結果を同期した。
\end{enumerate}

\subsection{実際の会話履歴}
% Paste the original conversation below.
\begin{Verbatim}
未添付：実際のAI会話履歴をここに貼り付けてください。
\end{Verbatim}
\end{document}
